\section{战争的最后二十年：中枢收缩与北方重组}

\subsection{战报之外的战略困局}

蒙古攻金之后，金廷面对的不是一条可由单次决战解决的战线。北方城镇相继受威胁，中都和河北的联系变得不可靠；南方仍有同南宋的边界、驻军和外交；中原的粮区、河道和城市又成为维持朝廷的最后支点。每一次把兵力送往一个方向，都会使另一个方向的道路、仓储或地方守备变薄。后世将这种局势概括为“腹背受敌”，虽抓住了压力，却容易遮住当时决策者面对的具体资源限制。

军事行动在材料中通常最醒目，原因是帝纪和列传以战事排列年代。我们需要把军队放回补给网络中理解。骑兵的机动需要草料和换马，步卒需要粮食和装备，围城需要木料、工匠、河运和长期驻留。蒙古军的作战方式、金军和地方守备的反应，都随着地形、季节和城防改变。不能把后来胜者的叙事当成战场全景，也不能用一方报称的斩获来替代对地方破坏的衡量。

\subsection{南迁不是逃离现实的终点}

宣宗向开封方向调整中枢，常被后人视为放弃北方的象征。更具体地看，这一决定承认中都已经难以同时承担宫廷、军政和供给的功能，也试图利用河南既有城市和粮区。可是，新中心需要从周边取粮、安置迁入者、恢复官署和维持河防；它把中都的危机带进了另一片早已承压的区域。迁都解决的是即时的中枢存续问题，代价是让河南和南方边界承担更多风险。

这一选择不应从君主品格推出。任何可选地点都有道路、仓储、地方关系和敌军距离的不同组合。留在中都可能失去朝廷，南迁可能失去北方联系；两种方案都不是无代价的。制度得失也因此格外清楚：集中官员和文书能维系国家名义，集中需求却会加重新都附近的征调。当地方生产与运输跟不上时，朝廷拥有的命令越多，能够兑现的承诺反而越少。

\subsection{宋蒙合作与其后的裂缝}

南宋同蒙古的接触和合作，不能被写成两方自然形成的同盟。南宋希望摆脱长期北方压力并取得河南旧地，蒙古则有自己的征服计划；双方在攻金时的共同目标有限，战后对城池、粮区和路线的理解很快不同。蔡州之役因此既是金亡的终局，也是宋蒙直接相邻的前奏。以后的冲突并非偶然背约，而是金这个中介消失后，原先被边界隔开的利益不得不在同一空间碰撞。

一二三四年的合作也让我们再次注意地方。军队协同行动需要各自的补给和向导，城破之后谁接收仓粮、安置居民、修复道路，都会改变短期联盟的意义。宋、蒙、金三方材料对责任和功劳的叙述并不相同，任何一方都可能将政治需要写成道德正当。读者可以确认蔡州陷落及金朝终止这一结果，却不能从结果倒推出每一支军队、每一座城和每一户居民都具有相同的目标。

\subsection{城陷之后，制度在哪里继续}

王朝终结并不会立即抹去它的制度。金代的道路、州县名称、税役分类、寺院和市场仍由留在当地的人使用、回避或重新解释；后来的蒙古政权也必须面对同样的粮食、交通和登记问题。金朝统治留下的遗产不是一套被完整继承的法典，而是一片被战争、迁徙和行政反复重塑的社会空间。后继者可以改换官名和税目，却不能凭一纸命令恢复被毁的桥梁、离散的家庭或失去的地方知识。

这种延续也是理解金亡为何不只是一个朝代故事的关键。若只以帝号存亡排列历史，蔡州之后像是章节的句号；若看道路、家户和地方机构，它更像一个漫长重组中的断点。有人把旧经验带入新政权，有人因旧身份而失去位置，有人只是在新的征调与新的市场中继续谋生。材料保存得越少，越应避免给他们安排整齐的政治台词。
